\documentclass[9pt,a4paper]{extarticle}
\usepackage[T1]{fontenc}
\usepackage[utf8]{inputenc}
\usepackage[spanish,es-noquoting]{babel}
\shorthandoff{"}
\usepackage{libertine}
\usepackage[scaled=0.80]{beramono}
\usepackage[margin=8mm]{geometry}
\usepackage{multicol}
\usepackage{xcolor}
\usepackage{tikz}
\usetikzlibrary{arrows.meta,positioning,fit,calc}
\definecolor{tinta}{HTML}{12191A}
\definecolor{acento}{HTML}{B8501C}
\definecolor{verde}{HTML}{146B5F}
\definecolor{gris}{HTML}{5A6A6B}
\definecolor{fondo}{HTML}{EDF1F0}
\definecolor{linea}{HTML}{C8D4D1}
\setlength{\columnsep}{5mm}
\pagestyle{empty}
\begin{document}
\color{tinta}
\begin{multicols}{2}
\begin{center}
\begin{tikzpicture}[
  font=\scriptsize,
  caja/.style={draw=linea,fill=white,rounded corners=1pt,inner sep=1.2mm,text width=30mm,align=center},
  agente/.style={caja,draw=acento,line width=0.6pt},
  herr/.style={caja,draw=verde,line width=0.6pt,text width=28mm},
  fl/.style={-{Stealth[length=1.4mm]},draw=gris},
  et/.style={font=\tiny,text=gris,align=left,inner sep=0.6pt},
  node distance=5.6mm]

\node[caja,fill=fondo] (q) {pregunta en prosa};
\node[agente,below=of q] (a) {\textbf{A · Intérprete}\\ lee la pregunta};
\node[herr,below=7mm of a] (t1) {\textbf{T1 · buscar\_ramo}\\ código, sin modelo};
\node[agente,below=of t1] (b) {\textbf{B · Resolvedor}\\ elige entre 5 candidatos};
\node[agente,below=7mm of b] (c) {\textbf{C · Verificador}\\ pregunta + calendario};
\node[herr,below=9mm of c] (t2) {\textbf{T2 · motor de reglas}\\ el código del D2};
\node[caja,fill=fondo,below=of t2] (out) {\textbf{fin} · decisión y regla};
\node[caja,draw=gris,dashed,text width=11mm,right=8mm of a] (fin) {\textbf{fin}\\ sin ramo};
\node[caja,fill=fondo,text width=11mm,left=7mm of t2] (hist) {historial};

\draw[fl] (q) -- node[right,et]{texto} (a);
\draw[fl] (a) -- node[right,et]{consultado (otros y período\\ siguen a B, C y T2)} (t1);
\draw[fl] (t1) -- node[right,et]{5 candidatos, marca\\ la coincidencia exacta} (b);
\draw[fl] (b) -- node[right,et]{elige o mantiene: código} (c);
\draw[fl] (c) -- node[right,et]{acepta, o corrige\\ el período: código, período} (t2);
\draw[fl] (t2) -- (out);
\draw[fl] (hist) -- (t2);
\draw[fl,dashed] (a) -- node[above,et]{sin\_ramo} (fin);

% la pregunta también llega al Verificador
\draw[fl,draw=linea] (q.east) -- ++(25mm,0) |- node[pos=0.3,left,et]{pregunta} (c.east);
% bucle del Resolvedor con T1
\draw[fl] (b.east) -- ++(4mm,0) |- node[pos=0.25,right,et]{busca de\\ nuevo (≤2)} (t1.east);
% revisión de sus citas
\draw[fl] ($(c.south)+(-6mm,0)$) .. controls +(0,-5mm) and +(0,-5mm) .. node[pos=0.5,below,et]{revisa sus citas (≤1)} ($(c.south)+(-13mm,0)$);
% tope
\draw[fl,dashed] ($(c.south east)+(0,0)$) -- ++(9mm,0) |- node[pos=0.25,right,et]{tope:\\ última\\ elección} (t2.east);

% vueltas al Intérprete y al Resolvedor
\draw[fl] (c.west) -- ++(-4mm,0) |- node[pos=0.25,left,et,align=right]{rechaza:\\ código\\ + motivo} (b.west);
\draw[fl] ($(c.west)+(0,-1.4mm)$) -- ++(-14mm,0) |- node[pos=0.2,left,et,align=right]{rechaza:\\ mención\\ + motivo} (a.west);
\draw[fl] ($(b.west)+(0,1.4mm)$) -- ++(-8mm,0) |- node[pos=0.3,left,et,align=right]{ninguno} ($(a.west)+(0,-1.4mm)$);
\end{tikzpicture}
\end{center}

{\footnotesize\color{gris} Borde naranjo, agentes, cada uno con su archivo en \texttt{prompts/}.
Borde verde, herramientas. Las vueltas por rechazo o \texttt{ninguno} suman máx. 2; el tercer
rechazo, o 12 llamadas, cierra el caso con la última elección.\par}

\columnbreak
\mbox{}
\end{multicols}
\end{document}
